\subsection{The elliptic pullback at logarithmic order two}
\label{subsec:elliptic-endpoint-geometry}

At the smallest characteristic-two boundary, the curve-pullback bound
leaves an elliptic base.  The two logarithmic forms on the distinguished
curve determine the exact inseparable degree of that multisection and
constrain the first normalized base change of the surface.

\begin{proposition}[Exact height and ramification]
\label{prop:elliptic-endpoint-height}
Let $X$ be a smooth integral projective surface over an algebraically
closed field $k$ of characteristic two, and let $C\subset X$ be a
smooth integral ample divisor with $d=C^2$.
Suppose $f:X\to E$ is a morphism to an ordinary elliptic curve and
$P\in\Pic^0(E)$ has exact order two, with
\[
 M=f^*P\ne\OO_X,\qquad M|_C\simeq\OO_C.
\]
Assume that the logarithmic form of $M$ has exact order two along $C$
and that $2g(C)-2=3d$.  Let $Z$ be the leading tangential divisor in
Proposition~\ref{prop:smallest-endpoint-rank}, and put
$n=\deg(f|_C)$.
The $P$-torsor $r:E_1\to E$ is a smooth integral degree-two purely
inseparable cover of genus one.  There is a factorization
\begin{equation}\label{eq:elliptic-endpoint-factorization}
 f|_C=r\circ h,\qquad h:C\longrightarrow E_1,
\end{equation}
where $h$ is separable of degree $n/2\geq2$.
Thus the inseparable degree of $f|_C$ is exactly two and $n\geq4$.
The ramification divisor $R_h$, with its different multiplicities,
satisfies
\begin{equation}\label{eq:elliptic-endpoint-ramification}
 \deg R_h=3d,\qquad
 \operatorname{Supp}R_h\cap\operatorname{Supp}Z=\varnothing.
\end{equation}
In particular $h$ is unramified at every point of
$\operatorname{Supp}Z$.
\end{proposition}

\begin{proof}
The nontriviality of $M$ makes $f$ nonconstant.
Proposition~\ref{prop:boundary-curve-pullback-excess} shows that $C$
dominates $E$ and that $n$ is even.  Multiplication by two on
$\Pic^0(E)(k)$ is surjective.  Choose $Q$ with $Q^{\otimes2}\simeq P$
and put $L=f^*Q$.  Then $L^{\otimes2}\simeq M$, so the hypotheses
of Proposition~\ref{prop:smallest-endpoint-rank} hold; in particular
the required N\'eron--Severi divisibility follows from this actual
square root.

The $P$-torsor is integral.  Indeed, a square root of one of its
defining units in $k(E)$ would extend to local unit roots by normality,
and the unique roots would trivialize $P$.  Its logarithmic form
$\omega_P$ is nonzero by the injectivity of the logarithmic-form map
on actual two-torsion of a smooth proper curve, as used in
Lemma~\ref{lem:positive-square-resolved-transfer}.
Since $E$ has genus one, $\omega_P$ is nowhere zero.
The local equation $z^2=u$ of the torsor is therefore smooth, since
$du/u=\omega_P$ has no zero.  Thus the torsor is already normal and
smooth.  Its function field is $k(E)^{1/2}$, so it is an inverse
Frobenius twist of $E$ and has genus one.  The map $r$ is finite,
flat and purely inseparable of degree two.

Set $Q'=r^*Q$.  Its square is trivial, since the torsor trivializes
$r^*P$.  It is nontrivial: otherwise the norm formula would give
\[
 \OO_E\simeq\operatorname{Nm}_r(r^*Q)
       \simeq Q^{\otimes2}\simeq P.
\]
Its logarithmic form $\theta$ is therefore a nonzero, nowhere-zero
regular differential on $E_1$.

The pullback of the $P$-torsor to $C$ is trivial.  Indeed,
$M|_C\simeq\OO_C$, and in a global frame the square trivialization
is a nonzero constant, which has a square root in $k$.
A section of this pulled-back torsor gives the morphism $h$ in
\eqref{eq:elliptic-endpoint-factorization}.
It is nonconstant, hence finite, and
\[
 U=L|_C=h^*Q',\qquad \beta_U=h^*\theta.
\]
By Proposition~\ref{prop:smallest-endpoint-rank}, the form $\beta_U$
is nonzero and is nonvanishing on $\operatorname{Supp}Z$.
A finite map between smooth curves is separable exactly when its
function-field differential is nonzero.  Thus $h$ is separable.
Degree comparison gives $\deg h=n/2$.
The integer $d=\deg Z$ is positive and even, so
$g(C)=1+3d/2\geq4$.
A finite degree-one map between smooth projective curves is an
isomorphism; consequently $h$ cannot have degree one.
The factorization has a separable part of degree $n/2$ and a purely
inseparable part of degree two, proving the asserted exact height.

For the separable map $h$, the differential map
$h^*\omega_{E_1}\to\omega_C$ has zero divisor $R_h$, with the
different multiplicities also in the presence of wild ramification.
Since $\theta$ trivializes $\omega_{E_1}$, this says
\[
 R_h=\operatorname{div}(h^*\theta)
    =\operatorname{div}(\beta_U).
\]
Its degree is $2g(C)-2=3d$, and its support avoids
$\operatorname{Supp}Z$.  This proves
\eqref{eq:elliptic-endpoint-ramification}.
\end{proof}

\begin{proposition}[Smoothness after the first normalized base change]
\label{prop:elliptic-endpoint-smoothing}
Under the hypotheses of Proposition~\ref{prop:elliptic-endpoint-height},
let $W$ be the normalization of $X\times_E E_1$, and write
\[
 \pi:W\longrightarrow X,\qquad
 \varphi:W\longrightarrow E_1,\qquad B=\pi^*C.
\]
Then $\varphi$ is smooth at every point of $B$.
Its non-smooth locus is finite, and its generic fibre is smooth.
If $f_*\OO_X=\OO_E$, its generic fibre is also geometrically integral.
\end{proposition}

\begin{proof}
The base change $X\times_E E_1$ is the $M$-torsor, so $W$ is its
normalized torsor.  Lemma~\ref{lem:positive-square-root-genus}
makes $B$ a reduced integral effective Cartier divisor, and the proof
of Corollary~\ref{cor:positive-excess-characteristic-two-boundary}
shows that $W$ is smooth on a neighbourhood of $B$.
Use $Q,L,Q'$ and $\theta$ from the preceding proof.
The pullback of $L$ to $W$ is $\varphi^*Q'$, and its logarithmic
form is
\[
 \eta=\varphi^*\theta.
\]
On the suitable resolution $S\to W$, its restriction to the smooth
strict transform $D$ is the nowhere-zero conormal section
\eqref{eq:smallest-endpoint-conormal}.

Every point $w\in B$ has a point of $D$ above it.
If the cotangent vector $\eta(w)$ vanished, its pullback at that
point would vanish in $\Omega_S^1|_D$, contradicting the nonzero
conormal section.  Thus $\eta$ is nonzero at every point of $B$.
The surface $W$ is smooth there, and $\theta$ generates the
cotangent line of the smooth curve $E_1$.
The differential criterion for smoothness therefore shows that
$\varphi$ is smooth along $B$.

The non-smooth locus of a finite-type morphism is closed.
The divisor $B$ is ample, since it is the finite pullback of the
ample divisor $C$.  A curve in the non-smooth locus would have to
meet $B$: the restriction of an ample line bundle to a complete
curve has positive degree.  This contradicts smoothness along $B$.
The non-smooth locus is therefore finite.  Its image in $E_1$ is
finite as well, so the generic fibre of $\varphi$ is smooth.

Finally assume $f_*\OO_X=\OO_E$, and put $F=k(E)$ and $K=k(X)$.
The Stein condition says that $F$ is algebraically closed in $K$.
Also $f^*\omega_P=\omega\ne0$, so the map
$\Omega^1_{F/k}\to\Omega^1_{K/k}$ is nonzero.
A separating parameter of $F/k$ consequently has nonzero
differential in $K$ and can be extended to a separating
transcendence basis of $K/k$.  It follows that $K/F$ is separably
generated.  Thus $K/F$ is a regular field extension and remains
geometrically integral after extension to $k(E_1)$.
Normalization does not change this generic function field.
The smooth generic fibre of $\varphi$ is therefore geometrically
integral under the Stein hypothesis.
\end{proof}

The smooth-generic-fibre assertion uses ampleness of $C$.
The separate fibration on the covering resolution given by
Proposition~\ref{prop:boundary-resolved-fibration} already follows
from positive self-intersection.  The simultaneous occurrence of
these structures remains a necessary condition on an elliptic
endpoint configuration.
